\section{A Anatomia do Game Loop}

\begin{frame}{O Ciclo de Vida de um Jogo}
    Todo jogo de computador, do clássico Pong aos títulos AAA modernos, é executado dentro de quatro fases essenciais a cada quadro (\textit{frame}):

    \vspace{0.3cm}
    \begin{center}
    \begin{tikzpicture}[
        node distance=1.0cm and 1.5cm,
        box/.style={rectangle, draw=blue!70, fill=blue!10, rounded corners, text centered, minimum width=2.8cm, minimum height=0.7cm, font=\footnotesize\bfseries},
        arrow/.style={-{Stealth[scale=1.0]}, thick, color=blue!80}
    ]
        \node[box] (eventos) {1. Entrada (Eventos)};
        \node[box, right=of eventos] (logica) {2. Lógica (Update)};
        \node[box, below=of logica] (desenho) {3. Desenho (Render)};
        \node[box, below=of eventos] (clock) {4. Relógio (Clock Tick)};

        \draw[arrow] (eventos) -- (logica);
        \draw[arrow] (logica) -- (desenho);
        \draw[arrow] (desenho) -- (clock);
        \draw[arrow] (clock) -- (eventos);
    \end{tikzpicture}
    \end{center}

    \begin{itemize}
        \item \textbf{1. Entrada:} Lê cliques do mouse, teclas pressionadas e fechamento da janela.
        \item \textbf{2. Lógica:} Atualiza posições, calcula colisões, pontuação e regras.
        \item \textbf{3. Desenho:} Limpa o buffer de vídeo e redesenha todos os elementos.
        \item \textbf{4. Relógio:} Aguarda o tempo necessário para travar a taxa de quadros (ex: 60 FPS).
    \end{itemize}
\end{frame}

\begin{frame}[fragile]{O Código Mínimo Essencial (Esqueleto)}
    \begin{lstlisting}[language=Python]
import pygame, sys

pygame.init()
tela = pygame.display.set_mode((800, 600))
pygame.display.set_caption("Meu Primeiro Jogo")
relogio = pygame.time.Clock()

executando = True
while executando:
    # 1. Eventos
    for evento in pygame.event.get():
        if evento.type == pygame.QUIT:
            executando = False

    # 2. Logica (vazio no momento)

    # 3. Desenho
    tela.fill((30, 30, 45))  # Limpa fundo com cinza escuro
    pygame.display.flip()    # Exibe no monitor

    # 4. Controle de FPS
    relogio.tick(60)

pygame.quit()
sys.exit()
    \end{lstlisting}
\end{frame}

\begin{frame}{Por que o programa trava sem \texttt{pygame.event.get()}?}
    \begin{block}{O Sistema Operacional e a Fila de Mensagens}
        \begin{itemize}
            \item Os sistemas operacionais (Linux, Windows, macOS) enviam mensagens contínuas para as janelas abertas (foco, teclado, redimensionamento, botão de fechar).
            \item Se o jogo não consumir essa fila através de \texttt{pygame.event.get()}, o sistema assume que o processo entrou em congelamento (\textit{loop infinito travado}) e exibe a mensagem de ``Não Respondendo''.
        \end{itemize}
    \end{block}

    \vspace{0.3cm}
    \begin{columns}
        \begin{column}{0.5\textwidth}
            \textbf{O que faz \texttt{pygame.QUIT}?}
            \begin{itemize}
                \item Representa o clique no botão vermelho de fechar janela `X'.
                \item Permite sair suavemente do laço \texttt{while}.
            \end{itemize}
        \end{column}
        \begin{column}{0.5\textwidth}
            \textbf{O que faz \texttt{display.flip()}?}
            \begin{itemize}
                \item Implementa o \textbf{Double Buffering}: tudo é desenhado na memória oculta e depois trocado de uma só vez para a tela visível, evitando cintilação (\textit{flickering}).
            \end{itemize}
        \end{column}
    \end{columns}
\end{frame}
